\section{The common history and its centered service}
\label{sec:selected-centered-service}
We now apply the local constructions to a single family. The input is
the simultaneous purchase of
Sections~\ref{sec:common-funded-history} and
\ref{sec:common-observable-purchases}, including the stationary and
free convergence estimates on their common horizon. We retain the
prescribed positive blocking penalty throughout.

Choose nondecreasing profile orders and spatial powers, both at least the
stage index, and nondecreasing support exponents at least $100$. The
particular choices $6+j$, $2+j$ and $100+j$ suffice for the elementary
construction. The selected public run uses the first stage $4$, the
orders $6+j$, the powers $2+j$, and the constant support exponent
$b_j=100$ at every stage; the identity $b_j=100$ is retained with that
run and is required by later consumers. Throughout this section the
public run $R$ is retained, including its selected shooting parameters,
its reference schedule $t_R$, its transition records and its transfer
family. The following paragraphs unpack the construction and the
estimates of these retained witnesses; where they speak of choosing an
object, the object meant is the one stored in $R$, and no second choice
is made. Choose the fixed parameter radius and the small response reserve
as in Section~\ref{sec:common-observable-purchases}. Make the block-length
choices before their numerical costs. Apply
Section~\ref{sec:common-scale-selection} to all the resulting requests,
with a nondecreasing occupancy schedule. This gives one sequence of
levels, couplings and periods, with $J(m)\to\infty$.

\subsection{Choosing the history before its observable comparisons}
For every cutoff $N$ and period $p$, the terminal family is the exact
Wilson family of Section~\ref{sec:wilson-startup-construction}.
The interval return and scalar reset are applied with the strengthened
local conditions of Section~\ref{sec:common-observable-purchases}.
Thus every member of the resulting continuous interval family carries
those conditions. To prove that its bottom couplings cover the required
interval, forget only these additional conditions and apply the barrier
and covering argument of Section~\ref{sec:surviving-interval-shooting}.
This leaves the interval, its endpoints and every replayed history
unchanged. Choose a parameter in that very interval whose bottom
coupling is the prescribed common value, and retain its original
strengthened certificate.
Retain that member, its actual representations $R_{N,p,m}$ and its
adjacent transitions. All of them use the already purchased paths,
kernels and clipping data. Choosing the members for the countably
many pairs $(N,p)$ introduces no new scale request. On the run $R$ these
are stored data: for every $(N,p)$ the run stores the interval family, a
parameter in its interval, and the identity of its history at $(N,p)$
with the replay of that family at that parameter. Its nesting receipt
records that every adjacent pair of this history is a protected and
accurate transition of the purchased step and that the terminal
representative carries the reserved addresses.

The original certificate at the chosen parameter supplies the local
parameter estimates on these histories. Thus the comparisons and fixed-level limits in
Sections~\ref{sec:coupled-background-convergence},
\ref{sec:selected-history-limits} and
\ref{sec:fixed-rung-convergence} apply to this family. The uniform
local derivative bank is that of
Section~\ref{sec:common-observable-purchases}. In particular, the
fixed-level convergence is available for the actual selected
activity arrays, not just for their stationary paths.

This application carries one input that the certificate of the
selected member does not by itself provide.
Section~\ref{sec:coupled-background-convergence} compares the selected
history with auxiliary truncated histories, which have a prescribed
coupling script and zero terminal activity, and
Section~\ref{sec:selected-history-limits} draws its limits under the
assumption that the transition hypotheses of
Section~\ref{sec:parameter-comparison-suppliers} hold for the actual
input of each such auxiliary history at every step, on the same path
and dictionaries. The certificate at the chosen parameter concerns
the selected history only. The fixed-level limits are therefore used
here subject to that same assumption on the auxiliary histories; it
is not proved in this section.

Put $t_d=t_R(d)$, the reference level stored in the run for the grade
$d$. The public construction supplies
\[
                 1\le t_d,\qquad d+3\le J(t_d).
\]
Such levels exist by divergence of $J$; no monotonicity of
$d\mapsto t_d$ is recorded or used. At each adjacent step
choose the wing containing the origin principal modes and the keeper
keys for all relevant support/inventory shapes. These keys form a finite
set: the torus is finite, so its support powerset and the pairs
$I\subseteq X$ with $I\ne\varnothing$ are finite. The wing uses
these shapes, not the countably many reserved serial addresses. Hence
it covers every forced coefficient while remaining finite. Take the
finite union required for the admitted grades.
Section~\ref{sec:heading-history-handoff} supplies one heading and
one local transfer map for that entire wing. Consequently the same
map works for all those grades. Enlarging the wing preserves its
finite interpolation proof and the exact endpoint representation.

Apply the finite backward construction of
Section~\ref{ex-centered-freight-producer}, starting at the terminal
marked section constructed in Section~\ref{sec:terminal-marked-construction}. This gives the finite transfer data at every grade,
cutoff and period. The three exact identities in
\eqref{eq:history-three-identities} identify the coefficient,
exponential entry and forced profile throughout the construction.
In particular, its terminal density is
$a_N^d R_{N,p,N}(U)\,\mathrm{initial}_{N,p}(v)(U)$;
the physical observable $\mathrm{initial}_{N,p}$ already contains
the factor $a_N^{-d}$. The local maps used
in this construction are the adjacent maps already chosen above;
the downward induction does not select another background history.
The unit, gauge invariance and signed-motion identities are preserved
by the exact local section and by the equivariant finite assembly.
Thus these transfer data form one family on the retained history. On the
run $R$ this family is the stored transfer family $S_R$, and the adjacent
maps are its stored transition records; for $t_d>N$ the stored transfer
datum is a default map and carries no assertion, as explained in the
proof of Proposition~\ref{prop:ex-original-history-attachment}.

\subsection{Paid retention on the stored doors}

Write $x_i=\widehat\beta_i$, $\alpha_i$ for the support exponent,
$B_i=B(x_i)$, and $r_i=N_{J(i+1)}$. Set
\[
 y_i=x_{i+1}^{-1},\qquad \delta_i=x_{i+1}^{-2},
 \qquad u_i=\tfrac14 x_i^{-1/2}.
\]
Use $w_i^{\mathrm f}$ for the history's forced allowance and $w_i$ for
the distinct flagged \emph{response} price of the tariff at step $i$.

\begin{proposition}[Retention and terminal endowment on the selected history]
\label{prop:selected-centered-retention}
Keep the selected replay, nesting and doors above, and the fixed
profile and purchased entrance delays. At each stored interval point
assume the arrival, ordinary and raised-response
data of Section~\ref{sec:finite-menu-history-handoff}
are available on the retained representations and door. Assume also
the following preservation property of the retained return: for each
nonterminal selected entry, its output dictionary has zero regular
addresses for nonempty supports and eligible forced addresses for
every nonempty inventory contained in its support. Eligibility here
is required whenever the actual flags realize the inventory. Assume the delays include the tariff thresholds of
Proposition~\ref{prop:ex-lattice-eventual-tariff} for these targets and
the thresholds of Theorem~\ref{thm:ex-selected-geyser-steam}.
Then for every $N,p$ and $0\le i<N$ the stored door has the following
paid response. Its input is any finite priced offer of nonempty modes
with positive marker weights, fine sites in the offered supports, and
parent-site occupancy at most $r_i$. The site-rooted input price is at
most $y_i$ at exponent $\alpha_{i+1}-4$. Its ordinary and flagged
output families have the support and flag conditions and per-mode
majorants in the multi-marker offer of
Section~\ref{sec:actual-marked-portfolio-calibration}.
Their rooted weighted prices have bounds $1$ and $w_i$ at exponents
$\alpha_i+1/2$ and $\alpha_i+1$, respectively. The assembled density is
the exact push of the offered polynomial. In each flagged box the
quiet coefficient is the lower ordinary activity profile and the
singleton coefficient is the sum of the corresponding offered
singleton profiles transported by this door. The door is sound and
tuned, and the upper contact addresses are retained. For the same $y_i,w_i$,
\[
 0<\delta_i\le1,\qquad \alpha_i,\alpha_{i+1}\ge100,\qquad
 T_i+e^{T_i}T_i\le y_i,\qquad
 (2B_i+1)^4c_i\le\tfrac12,\qquad e^{c_i}-1\le1,
\]
where $T_i=r_i\mathsf s(B_{i+1})\delta_i+4\delta_i$,
$c_i=w_i+r_i\mathsf s(B_i)\delta_i$ and
$\mathsf s(B)=e\{2(16B+9)\}^4$.

At every terminal level $N\ge0$ the same schedule gives
\[
 u_N\ge0,\qquad \alpha_N\ge3,\qquad
 (2B_N+1)^4 C_N^{\rm end}\le3,\qquad e^{C_N^{\rm end}}-1\le1,
 \qquad
 C_N^{\rm end}=w_N^{\mathrm f}+r_N\mathsf s(B_N)u_N.
\]
\end{proposition}

\begin{proof}
For fixed $N,p,i$, take the interval witness of the stored step.
Its upper representation is the interval input at the retained point
and its lower representation the return at that same point. The door
is the one-step transport at that point, with those two endpoint
identities substituted. Thus replacing the local representations by
the selected entries in the response formulas changes neither the
operator on activities nor the dictionary. At the top,
Proposition~\ref{prop:wilson-startup-family} supplies the regular and
forced addresses on the selected Wilson representative. At any
lower entry the preservation property applies to the return of the
selected step above it. Hence the upper entry of every door has these
addresses, including when that entry is terminal. For a nonempty
support take its zero regular address. For each nonempty inventory
take an eligible forced address at that support, using eligibility
on the configurations whose flags realize the inventory.

At this step the entrance stage is $J(i+1)$, the target stage is
$J(i)$, and the coupling is $x_i$. The selected clock gives
$x_{i+1}=x_i+c_G\log L_i$ and the purchased entrance bound at $x_i$.
The tariff therefore applies at the retained occupancy $r_i$.
Its shipment inequality converts the site-rooted offer at the upper
exponent into the rooted birth-region price, using the parent map
and the occupancy bound, as in
Theorem~\ref{thm:ex-lattice-volant-supplier}. The ordinary response uses
the ordinary component budget and the raised response uses the
separately paid flagged budget. The latter has price $w_i$ from the
same tariff. Applying the finite-menu handoff to the retained
response data gives the two grouped ledgers. For a fixed mode, sum the
configuration-independent absolute majorants over entries of its output
support or box. The triangle inequality bounds its grouped coefficient
by that sum. Since the mode weights are nonnegative and unchanged by
grouping, the weighted sum of these majorants is bounded by the
corresponding grouped support price. The support restrictions,
quiet and singleton identities, and exact push are the grouping
identities of Theorem~\ref{thm:ex-lattice-volant-supplier}. The finite-menu
handoff gives soundness and compatibility with grouped sums for this
same one-leaf transport. This uses the purchased response data for
every priced offer; it does not select a new door when the offer changes.
The probe, purchase, lodging and credit inequalities of the tariff
are precisely the displayed inequalities for $\delta_i,T_i,c_i$
after the fine-coupling substitution. In particular the allowance used
to buy $T_i+e^{T_i}T_i$ is the input price of the very same response.

For the terminal assertion, apply
\eqref{eq:ex-geyser-second} at $x=x_N$ with stages $J(N)$ and $J(N+1)$.
This is a numerical assertion about the global schedule, so it uses
the purchased next entrance inequality even when the finite history
itself stops at $N$. The selected allowance is
$w_N^{\mathrm f}=\operatorname{stain}(x_N)$ and the other term in that
bound is $N_{J(N+1)}\mathsf s(B_N)s(x_N)$, which equals
$r_N\mathsf s(B_N)u_N$. The resulting two inequalities give the claimed
endowment. Positivity of $x_N$ and $\alpha_N\ge100$ give the remaining
two assertions. This argument precedes the choice of a terminal
label and needs no terminal voice, including at $N=0$.
\end{proof}

The proposition concerns the stored doors of the public run, including
its step at level zero when $N>0$. The per-campaign service projection
uses the campaign's positive-level doors. Its retention assertion
is restricted to $1\le k\le i<N$; no response at its default zero-level
door is required. In either application the terminal endowment depends
only on the schedule. The preservation of the required slots by the
retained return is proved in Lemma~\ref{lem:selected-reserved-slots}
below. The remaining analytic attachment in the proposition is the
availability, on the selected replay, of the arrival and
raised-response receipts.

\begin{lemma}[Reserved slots on every selected entry]
\label{lem:selected-reserved-slots}
Keep the selected history. For every $N,p$ and every $0\le m\le N$,
the stored representation $R_{N,p,m}$ is written in the fixed terminal
dictionary of Proposition~\ref{prop:wilson-startup-family}, and its
coefficients on the reserved labels vanish. Precisely:
\begin{enumerate}
\item for every regular address and every natural number, the regular
  coefficient of $R_{N,p,m}$ on the corresponding reserved regular
  label is zero;
\item for every support $X$, every nonempty inventory $I\subseteq X$
  and every natural number, the forced coefficient of $R_{N,p,m}$ on the
  corresponding reserved forced label is zero. That label has support
  $X$ and inventory $I$, and its additional predicate holds at every
  configuration $U$ with $X\cap\operatorname{flags}(U)=I$.
\end{enumerate}
In particular, the preservation property assumed in
Proposition~\ref{prop:selected-centered-retention} holds for every
nonterminal selected entry. Every stored representation, terminal or
not, has a Keeper, that is, an injective eligible assignment of a
reserved forced address to every pair $\varnothing\ne I\subseteq X$.
\end{lemma}

\begin{proof}
For $m=N$, the stored top representation is a member $F_\beta$ of the
exact startup family, by the choice of the history at the beginning of
this section. Proposition~\ref{prop:wilson-startup-family} gives (1)
and (2) for $F_\beta$.

Let $m<N$. By the stored identities recalled in the proof of
Proposition~\ref{prop:selected-centered-retention}, $R_{N,p,m}$ is the
output of the interval return of step $m$ at the retained interval
point. The return produces its output in the same fixed terminal
dictionary as the startup. It reorganizes its renewed regular and
forced arrays by the same fixed packing and grouping; see the first
subsection of Section~\ref{sec:terminal-marked-construction}, where
this is stated for the startup and every return. The proof of
Proposition~\ref{prop:wilson-startup-family} shows that the vanishing
on reserved labels is a property of this packing and grouping. It holds
for every dictionary and every coefficient family to which they are
applied. The packing injection takes values among the unreserved
labels, and every entry outside its image is zero. A reserved forced
label is the group formed by one reserved hard label alone. Every
grouped label of the packed family consists of packed labels, so it is
not a reserved forced label, and every grouped entry outside the image
of the packed family is zero. Applying this to the renewed arrays of
the return gives (1) and (2) for $R_{N,p,m}$.

The support, inventory and additional predicate of a reserved label
are data of the dictionary, not of the coefficients. They are therefore
the same as for the startup family. Their predicate is the flag
condition, by the construction recalled in the first subsection of
Section~\ref{sec:terminal-marked-construction}.

For the preservation property, take any nonempty support. Its reserved
regular label is a regular address of $R_{N,p,m}$ with zero
coefficient. For a nonempty inventory $I\subseteq X$, its reserved
forced label is a forced address with support $X$ and inventory $I$.
That address is eligible whenever the actual flags in $X$ realize $I$.
Choosing, for each pair $(X,I)$, the reserved forced label with serial
number zero gives an injective assignment. This is the asserted Keeper.
\end{proof}

\subsection{The service at every reference level}
Fix a grade $d$, a period $p$, a cutoff $N$, and
$t_d\le k\le N$; in particular $k\ge1$. We verify the inputs of
Theorem~\ref{thm:ex-selected-run-service} for this family.

Proposition~\ref{prop:selected-centered-retention} supplies the adjacent
retention and terminal endowment under its stated selected-replay
hypotheses. Separately,
Proposition~\ref{prop:terminal-marked-voice} constructs the marked density
and its terminal forced-profile identity on the startup representation
when the reserved addresses, lawfulness, sound section and paid
entrance fee are available. The numerical endowment does not depend
on that label construction.

The room and bank constants are those chosen for the local
coefficient maps. The scale inequalities in
Section~\ref{ex-centered-geyser-scale-estimates} give their common
marker prices
\[
 \zeta_{d,m}=\widehat\beta_m^{-(\deg F_d+4)},\qquad
 \phi_m=\widehat\beta_{m+1}^{-1},\qquad
 \delta_m=\widehat\beta_{m+1}^{-2},\qquad
 \eta_m=\tfrac14\widehat\beta_m^{-1/2}.
\]
Here $F_d$ is the finite coefficient frame. Their strict positivity,
local coefficient bounds, adjacent-price growth and support inequalities
are proved there. Since $d\le J(k)\le J(m)$, the finite-grade
inequalities purchased at stage $J(m)$ apply at every $m\ge k$.
The occupancy bound at level $m$ is $N_{J(m+1)}$, which is
nondecreasing because both maps are nondecreasing.

The same scale estimates give
\[
 B_m\ge8,\qquad
 \delta_m\phi_m e^{B_m/128}\ge1,\qquad
 B_{m+1}^2\le L_mB_m^2.
\]
The stain/contact bound pays the simultaneous forced support
requirement. These inequalities continue to hold after restriction
to any reference level $k$; no constant is selected from $N$.
For the local-section norm use its fixed camera constant $C_d$.
The fitted stencil has cover at most one and its values are bounded
by $C_d$, uniformly in the cutoff and period. Its paid coefficient
price is precisely the bank price in the scale estimates. The
cover bound just stated is the footprint-cover clause of the height
record, item~(iii) of Theorem~\ref{thm:ex-selected-run-service}; the
support-size inequalities of its item~(ii) are among the scale
estimates displayed above. Section~\ref{ex-centered-site-geometry}
is not used for either: it supplies the occupancy, migration and
separation of the label sites, which enter only in the moment
bound below.

The backward transfer data supply the remaining finite transfer
record. Its coefficient and activity bounds are the two separate
ledgers in Section~\ref{ex-centered-freight-producer}; its adjacent
identities use the selected maps, and its terminal forced-profile
identity is the third identity in
\eqref{eq:history-three-identities}. That identity is essential:
the scalar integral identity alone would not give the terminal
marked polynomial. All these conclusions hold for every $(d,k,N,p)$
in the stated range, with the already fixed reference levels.
These are the remaining service fields under the stated local purchases,
finite backward construction and selected-replay hypotheses.

We may therefore apply Theorem~\ref{thm:ex-selected-run-service}
and the terminal birth calculation. At a reference level $k$, a
finite label list with paid weights and separated origins has
normalized moment bound
\[
              \exp\{|\Lambda_k|\mathcal T_k+6|\mathcal V|\}.
\]
Let $m\le k$ now denote the root level opened at $k$, as in
Section~\ref{ex-centered-site-geometry}, rather than a running level.
Take one label at each site of $\Lambda_m$, so that
$M=|\mathcal V|=|\Lambda_m|$, placed by the aligned embeddings of
Section~\ref{ex-centered-site-geometry}. With $n_j$ the side of the
periodic lattice $\Lambda_j$ and $q_k=a_m/a_k$, that section records
$n_k=q_kn_m$; since $|\Lambda_j|=n_j^4$, this gives the site identity
$|\Lambda_k|=M(a_m/a_k)^4$.
The normalized $M$th root is consequently bounded by
\[
              \exp\{(a_m/a_k)^4\mathcal T_k+6\}/w.
\]
This is the centered estimate for the same family, simultaneously
at all suitable periods. The fixed-level convergence just retained
and the removal of higher marker terms in
Section~\ref{sec:fullmesh-correlation-cauchy} then give convergence
along the full mesh sequence at every fixed suitable period.
The finite-difference bounds of
Section~\ref{sec:ex-fullmesh-fixedbounds} apply to the same
transfer data, so its limiting correlations have the smoothness
proved in Section~\ref{ex-smooth-limits}.

The further calibration of an arbitrary marked section and the
rank and ultraviolet comparisons must still be made on this family.
They will use the same reference levels and transfers; the centered
construction above is already fixed before those comparisons or a
volume subsequence are chosen.

\paragraph{Relation to the paid selected-history record.}
Section~\ref{sec:actual-marked-portfolio-calibration} consumes a record
(R1)--(R7) on the same family and the same transition records. Its
parameters are those of this section: the stage function $W$ is the
purchased worth, $u_i=\eta_i$, $\phi_i$, $\delta_i$ and $\zeta_{d,i}$
are the marker prices displayed above, $r_i=N_{J(i+1)}$, the doors are
the stored transition records, and $\psi_i$ is the common waybill
\eqref{eq:history-common-waybill} at level $i+1$. With these
parameters, the present section assembles the following items, each from
the inputs named in its place above. The room and bank inequalities (R1)
are those of the purchased budget. The scale bounds (R2) are the
positivity of the marker prices and the three displayed inequalities.
The monotonicity in (R3) is that of the occupancy bound. Proposition~\ref{prop:selected-centered-retention} gives (R4) under
its stated hypotheses. Of (R5), this section supplies $\phi_i\le1$,
$\psi_i\ge0$, the second, third and fourth displayed inequalities, the
bound $X=C_d$, the cover number one, and the terminal inequality for
$u_N$. The backward transfer data, with their door-by-door comparison at
price $\psi_i$ and their terminal forced-profile identity, are (R6).
Proposition~\ref{prop:selected-terminal-voice} below identifies the
terminal voice of (R6) on the stored top representation, with its
singleton hook. Proposition~\ref{prop:selected-music} attaches the
adjacent ledgers and finite-wing comparisons to one Music on that voice;
its one-step input remains open.
Three items are not part of the centered service and are not derived in
this section: $r_i\ge1$ in (R3); the positivity $\zeta_{d,i}>0$ and the
first inequality of (R5), which is
\eqref{eq:ex-marked-height-jump} for the prices above; and the paid
one-step input (R7). They belong to the stronger form of the funded
continuation described in the proof of
Proposition~\ref{prop:ex-original-history-attachment}: $r_i\ge1$ is the
prescribed lower bound of the occupancy schedule, positivity is that of
$\widehat\beta_i$, the first inequality of (R5) is the tail clause of
the delay purchased in Section~\ref{sec:common-observable-purchases},
and (R7) is the local transition estimate on the stored doors. For
(R7), Proposition~\ref{prop:selected-zero-coordinate-coupon} below
proves the tuning inequality from the grade-zero floor of
Section~\ref{sec:common-observable-purchases}, together with the signs.
It reduces the zero-coordinate Coupon to the regular and forced
insertion estimates, soundness of the stored door, and the countable
Coupon merger; these three inputs remain open on the selected history.
The proof of Proposition~\ref{prop:ex-original-history-attachment}
names this estimate but does not derive it. For none of the items has
this section a self-contained proof: each rests on the cited
producers.

\subsection{The stored terminal voice and its Music}
\label{sec:selected-terminal-voice}

This subsection treats item (R6) of
Section~\ref{sec:actual-marked-portfolio-calibration}, and the
corresponding clause of the service record consumed in
Theorem~\ref{thm:ex-selected-run-service}. Two things are needed. The
top of the backward transfer data must be the marked voice of
Proposition~\ref{prop:terminal-marked-voice} on the \emph{stored} top
representation $R_{N,p,N}$. Every adjacent comparison must then be
attached to one and the same family of voices, whose plays are the
stored transfer entries. The profile-difference identity in
\eqref{eq:history-three-identities} concerns one transition and an
arbitrary response $\psi$. It does not identify the terminal forced
profile. That identification is assertion~(4) of
Proposition~\ref{prop:terminal-marked-voice}, which we now place on the
stored representation.

Fix a grade $d$, a period $p$ and a cutoff $N$ with $t_d\le N$. Then
$d+3\le J(t_d)\le J(N)$. Write $R^\top=R_{N,p,N}$ for the stored top
representation, $S$ for the grade-$d$ local section with its fixed camera
constant $C_d\ge1$, and $w^{\mathrm f}_N$ for the forced allowance of
level $N$. A \emph{voice} at level $m$ is a complex-linear family
$v\mapsto V(v)$ of coefficient fields on $R_{N,p,m}$, together with a
linear coefficient index $x(v)\in F$ and a charge $c(v)\ge0$. We say that
a voice is \emph{heard} at $m$ if all of the following hold for every
$v$:
\begin{itemize}
\item $V(v)$ has the terminal format of
  Section~\ref{sec:selected-zero-coordinate-coupon} at the origin of
  $R_{N,p,m}$, and remainder bound $c(v)$ at weight
  $\sqrt{w^{\mathrm f}_m}$;
\item its coordinate vector is $x(v)$;
\item $c(v)\le s_m\|\operatorname{core}x(v)\|_1$, where $s_m\ge1$ is the
  column count at level $m$;
\item $c(u+v)\le c(u)+c(v)$;
\item the index commutes with the signed lattice motions.
\end{itemize}
A voice satisfies the \emph{profile covariance} at $m$ if, for every
motion $g$, transporting the forced profile of $V(v)$ by $g$ gives the
forced profile of $V(g\cdot v)$. A \emph{Music} is a choice of one voice
at every level $m\le N$.

\begin{proposition}[The stored terminal voice]
\label{prop:selected-terminal-voice}
Keep the selected history and the purchases of
Section~\ref{sec:common-observable-purchases}. Let $d,p,N$ be as above.
Then the four hypotheses of Proposition~\ref{prop:terminal-marked-voice}
hold for $R^\top$, its local section $S$ with constant $C_d$, and
$w=w^{\mathrm f}_N$. The marked voice
$\mathcal V^\top(v)=\mathcal V(v)$ of that proposition is therefore
defined on the stored top representation. Its index is the identity,
$x^\top(v)=v$, and its charge is $c^\top(v)=\|\operatorname{core}v\|_1$.
For every $v\in F$:
\begin{enumerate}
\item $\mathcal V^\top(v)$ has lead $v$ and zero regular part;
\item the density identity (T5) holds,
  $\operatorname{dens}\mathcal V^\top(v)=a_N^dR^\top\,
  \mathrm{initial}_{N,p}(v)$;
\item the forced profile of $\mathcal V^\top(v)$ on the activities of
  $R^\top$ equals the attachment of the singleton response $f_v$ on the
  footprint $Y$ (the singleton hook);
\item $\mathcal V^\top$ is heard at $N$ and satisfies the profile
  covariance at $N$.
\end{enumerate}
\end{proposition}

\begin{proof}
\emph{The stored top representation.} By the choice of the history at
the beginning of this section, $R^\top$ is the member $F_\beta$ of the
exact Wilson startup family of
Proposition~\ref{prop:wilson-startup-family}, at the coupling of the
retained parameter. That proposition applies at level $N$ under the
purchased schedule, and $F_\beta$ satisfies the retained-history
conditions at level $N$.

\emph{Hypothesis (1).} The forced arrays of $F_\beta$ have the locality,
measurability and gauge invariance required of the dictionary.
Proposition~\ref{prop:wilson-startup-family} bounds their forced cost by
the forced allowance $\delta_K$ of level $N$. The forced cost is the
rooted support sum
$\sup_x\sum_{i:\,x\in X_i}e^{E\operatorname{cov}_B(X_i)}k_i$ of a
majorant family $k_i\ge|K_i|$, with $E$ the support exponent; see the
packet majorant in the proof there. This is (T3) with $w=\delta_K$.
Moreover $\delta_K=\exp\{-M^{1/5}\}$ with
$M=\exp((\log\log\widehat\beta_N)^2)$, which is the allowance
$w(\widehat\beta_N)$ of Section~\ref{sec:terminal-marked-construction}.
This is the forced allowance $w^{\mathrm f}_N$ of the history.

\emph{Hypothesis (2).} Proposition~\ref{prop:wilson-startup-family}
states that the fixed terminal dictionary contains, for every forced
support $X$ with nonempty inventory $I$ and every natural number, a
reserved forced label with support $X$ and inventory $I$. It also states
that the forced coefficient of $F_\beta$ on every reserved label is zero.
The additional predicate of such a label is the flag condition
$X\cap\operatorname{flags}(U)=I$, by the construction of the dictionary
recalled in the first subsection of
Section~\ref{sec:terminal-marked-construction}. Hence hypothesis~(2)
holds. Selecting one reserved label for each pair $(X,I)$ gives an
injective assignment of an address with tag $(X,I)$, eligible whenever
$X\cap\operatorname{flags}(U)=I$. This is the Keeper of $R^\top$, and it
is the property recorded in the nesting receipt of the run.

\emph{Hypothesis (3).} The grade-$d$ local section is sound, and its
axis values are bounded by the fixed camera constant $C_d\ge1$ at every
configuration, as recorded in the service at every reference level
above.

\emph{Hypothesis (4).} Since $d\le J(N)$, the final sentence of
Proposition~\ref{prop:terminal-marked-voice} applies at the admitted
terminal level $N$. The fee thresholds are among the entrance delays of
Section~\ref{sec:common-observable-purchases}, as stated in the third
subsection of Section~\ref{sec:terminal-marked-construction}.

\emph{Conclusions.} Assertions (1)--(3) are the definition of
$\mathcal V(v)$ together with assertions (3) and (4) of
Proposition~\ref{prop:terminal-marked-voice}. For (4), the format
clauses are assertion~(1) there together with the zero regular part.
The origin lies in every footprint support $T_b\supseteq Y$. The
remainder bound with $a_{\rm r}=0$ and $a_{\rm f}=\|\operatorname{core}v\|_1$
at weight $\sqrt{w^{\mathrm f}_N}$ is assertion~(2) there. Indeed
$\operatorname{mass}(T_b,e)=M(T_b)e$, so the forced mass sum is at most
$\sqrt{w^{\mathrm f}_N}\,a_{\rm f}$. The cone bound is
$\|\operatorname{core}v\|_1\le s_N\|\operatorname{core}v\|_1$, since
$s_N\ge1$. Subadditivity is the triangle inequality for the core norm.
Covariance of the index is linearity of the motion action. Profile
covariance at $N$ follows from assertion~(5) of
Proposition~\ref{prop:terminal-marked-voice}, applied to the forced
profile through assertion~(4) there. The singleton response satisfies
$f_{g\cdot v}=g\cdot f_v$, and its attachment on $Y$ is transported by
the motion together with the footprint.
\end{proof}

For the adjacent steps we use the following one-step input. It is the
quantitative content of a step of Lemma~\ref{ex-freight-backward} on the
stored door, in the form in which the finite-wing comparison is exported.
For $m<N$, let $s_{m+1}$ be the column count at level $m+1$, $L_m$ the
block side and $\psi_m$ the common waybill
\eqref{eq:history-common-waybill}. A \emph{Carol} across the stored door
$D_m$ between voices $V'$ at $m+1$ and $V$ at $m$ consists of two
things. The first is a finite wing $W$. The second is, for every $v$, a
Tale witness on $W$ across $D_m$, in the sense of
Section~\ref{sec:selected-zero-coordinate-coupon}. This witness relates
the upper input data of $V'(v)$ to the lower forced profile of $V(v)$:
\begin{itemize}
\item the exponential entry and the forced profile on $R_{N,p,m+1}$;
\item the unit coordinate of $x'(v)$;
\item the physical factor $L_m^d$;
\item input price
  $C_{\rm orb}\|x'(v)\|_1\,(\psi_m+2s_{m+1})$ and output price $c(v)$.
\end{itemize}

\begin{enumerate}
\item[(RS)] \emph{One-step reprise on the stored door.} Let $m$ satisfy
  $t_d\le m<N$. Suppose $V'$ is a voice at $m+1$ that is heard and
  satisfies the profile covariance there. Then there is a voice $V$ at
  $m$ with four properties:
  \begin{itemize}
  \item $V$ is heard at $m$ and satisfies the profile covariance at $m$;
  \item for every $v$, $\operatorname{dens}V(v)$ is the push of
    $\operatorname{dens}V'(v)$ through the stored transition of step
    $m$, as in Lemma~\ref{ex-freight-backward};
  \item the coefficient indices and charges satisfy the paid adjacent
    cone and charge inequalities of
    Section~\ref{ex-centered-freight-producer};
  \item there is a Carol across $D_m$ between $V'$ and $V$.
  \end{itemize}
\end{enumerate}

The input (RS) is what the paid ticket of the stored door is to supply.
That ticket consists of:
\begin{itemize}
\item a finite wing containing the Keeper keys of $R_{N,p,m+1}$;
\item a heading;
\item the estimates of Section~\ref{sec:three-port-history};
\item soundness of $D_m$;
\item $\psi_m\ge0$;
\item $\sqrt{w^{\mathrm f}_{m+1}}\le1$;
\item the room and bank bounds.
\end{itemize}
It is not derived in this monograph on the selected replay. The
regular and forced estimates on the stored door are the input (K) of
Proposition~\ref{prop:selected-zero-coordinate-coupon}, and the slot
realization at nonterminal levels is the open attachment of
Proposition~\ref{prop:selected-centered-retention}.

\begin{proposition}[One Music on the stored terminal voice]
\label{prop:selected-music}
Let $d,p,N$ be as above, and assume (RS) for every $m\in[t_d,N)$. Then
there is a Music $(V_m)_{t_d\le m\le N}$ with the following properties:
\begin{enumerate}
\item $V_N=\mathcal V^\top$, the voice of
  Proposition~\ref{prop:selected-terminal-voice};
\item every $V_m$ is heard at $m$ and satisfies the profile covariance at
  $m$;
\item at every stored door $D_m$, $t_d\le m<N$, the density of $V_m(v)$
  is the push of that of $V_{m+1}(v)$, the adjacent cone and charge
  inequalities hold, and there is a Carol between $V_{m+1}$ and $V_m$
  with input price $C_{\rm orb}\|x_{m+1}(v)\|_1(\psi_m+2s_{m+1})$.
\end{enumerate}
The finite transfer data of Lemma~\ref{ex-freight-backward} may be taken
to be the plays of this Music, normalized by its reference readout. In
particular their terminal forced profile is the singleton hook of
Proposition~\ref{prop:selected-terminal-voice}.
\end{proposition}

\begin{proof}
We use downward induction on $k$ from $N$ to $t_d$. The claim at stage
$k$ is that there is a Music with $V_N=\mathcal V^\top$ such that (2)
holds at every $m\in[k,N]$ and (3) holds at every $m\in[k,N)$.

For $k=N$, put $V_N=\mathcal V^\top$, and fill the levels below $N$
with the zero voice. Assertion~(4) of
Proposition~\ref{prop:selected-terminal-voice} gives (2) at $N$; (3) is
vacuous.

Suppose the claim holds at stage $k+1$, with $t_d\le k<N$, and let
$(V_m)$ be the Music it provides. Its voice $V_{k+1}$ is heard and
satisfies the profile covariance at $k+1$. Apply (RS) at $m=k$ with
$V'=V_{k+1}$, and replace the voice at level $k$ by the resulting $V$.
The voices at levels above $k$ are unchanged, so (1), and (2) and (3)
above $k$, persist. At level $k$, (2) holds and the door $D_k$
satisfies (3) by (RS). This is the claim at stage $k$. After $N-t_d$
steps we obtain the Music.

The construction is finite. Each step modifies one level, and only
after the levels above it have been fixed; no voice is chosen after
testing a lower one. The density identities, the cone inequalities and
the terminal isomorphism $x_N=\mathrm{id}$ are those used in
Lemma~\ref{ex-freight-backward} and its normalization. The plays of
this Music are therefore a valid choice of the finite transfer data,
and their top is $\mathcal V^\top$.
\end{proof}

\paragraph{What these propositions supply.}
Proposition~\ref{prop:selected-terminal-voice} is unconditional relative
to the cited results. It discharges, on the stored top representation,
the four hypotheses of Proposition~\ref{prop:terminal-marked-voice}:
lawfulness, reserved addresses, the sound section and the paid fee. It
exports the lead, the zero regular part, the core charge, (T5), the
heard clauses and the singleton hook. The reserved addresses are thus
not an independent receipt of the run. They follow from the reserved
labels of the exact startup family. Proposition~\ref{prop:selected-music}
reduces the remaining part of (R6) to the one-step input (RS) on each
stored door. It places all the adjacent ledgers and finite-wing
comparisons on one Music whose top is that voice. It does not prove
(RS).

\subsection{The zero-coordinate Coupon on the stored doors}
\label{sec:selected-zero-coordinate-coupon}

This subsection treats item (R7) of
Section~\ref{sec:actual-marked-portfolio-calibration}. It proves the
tuning inequality from the purchase and the price bookkeeping
outright. The Coupon itself is reduced to three inputs on the same
stored door: the regular and forced insertion estimates, soundness,
and the countable merger of Coupons.

Fix a grade $d$, a cutoff $N$, a period $p$ and a level $i$ with
$t_d\le i<N$. Then $d+3\le J(t_d)\le J(i)$. Write
$R^+=R_{N,p,i+1}$ and $R^-=R_{N,p,i}$ for the stored upper and lower
representations, and $D=D_{N,p,i}$ for the stored door between them.
The block side of this step is $L_i$, the block length chosen at the
entrance stage $J(i+1)$. Write $C_i=\operatorname{columns}(i)$ for the
column count of the profile at the stage $J(i)$, and
$\kappa=\sqrt{w^{\mathrm f}_{i+1}}$, so that $0<\kappa\le1$. Put
\[
 \varrho_i=\operatorname{rate}(L_i)=L_i^{-2/5},\qquad
 \operatorname{barunit}(L_i)=L_i^{-3/8},
\]
and let $\psi_i=2e^{\alpha_{i+1}}\tau_{i+1}$ be the common waybill
\eqref{eq:history-common-waybill} at the fine level $i+1$; here
$\tau_{i+1}\ge0$ is the fine local-section derivative bound. The number
$\varrho_i$ is the number called $\eta_i$ in (R7). It is not the marker
price $\eta_m=\tfrac14\widehat\beta_m^{-1/2}$ displayed earlier in this
section, which is the allowance $u_m$ of (R5). Let $W_i$ be the finite
wing chosen above for this step. It contains the principal-mode keys of
the supports containing the origin and the keeper key of every
support--inventory tag.

\paragraph{Coefficients and their densities.}
We recall the data entering a Coupon. A coefficient field $Q$ on $R^+$
consists of a coordinate vector $Q.x\in F$, regular atoms
$(X_j,f_j)_{j\ge0}$ with supports $X_j$ and local functions $f_j$, and
a forced array $(k_a)_{a\ge0}$ indexed by the addresses of the
dictionary of $R^+$. Address $a$ carries a support $S_a$, a nonempty
inventory $I_a\subseteq S_a$ and an activity $K_a$; its \emph{tag} is
$\operatorname{tag}(a)=(S_a,I_a)$. The exponential entry of $Q$ is
\[
 \mathcal E(Q)(U)=Q.x_0+\ell(U,Q.x)
   +\sum_j\chi_{X_j}f_j(U),
\]
where $Q.x_0$ is the unit coordinate, $\ell(U,\cdot)$ is the masked
placement of $\operatorname{core}Q.x$ by the local section at the
origin, and $\chi_{X_j}f_j$ is the masked regular term. The forced
profile $\mathcal P(Q)$ is the function on tags
\[
 \mathcal P(Q)(X,I)=\sum_{a:\ \operatorname{tag}(a)=(X,I)}k_a .
\]
Call a finite set $b$ of addresses admissible at $U$ if it satisfies the
admissibility condition of the activity expansion of $R^+$. That
condition depends on the addresses only through their supports and
inventories, and two addresses of one admissible set have disjoint
supports. The dressed field at an address $a$ is
\[
 \operatorname{Dr}_a(U)=\sum_{b\ni a\ \text{admissible at }U}
            \ \prod_{c\in b\setminus\{a\}}K_c(U).
\]
The density of $Q$ is
\[
 \operatorname{dens}(Q)(U)=R^+.\mathrm{eval}(U)\,\mathcal E(Q)(U)
   +e^{R^+.\mathrm{exponent}(U)}\sum_a k_a(U)\operatorname{Dr}_a(U),
\]
in which the forced series has been regrouped by its address. The
regrouping is justified by the absolute bounds on the dressed sums
of a lawful representation and by the forced majorants below.

The terminal format at the origin requires several things of $Q$. Each
regular atom is regular on its support, its support contains the
origin, it vanishes at the identity configuration, and it passes the
holonomy test. Each forced coefficient is gauge invariant, measurable
and local to its support. Where $k_a$ is nonzero, address $a$ is
eligible, that is $S_a\cap\operatorname{flags}(U)=I_a$; and $k_a\ne0$
forces the origin to lie in $S_a$. A remainder bound $q$ at weight
$\kappa$ means that there are $a_{\rm r},a_{\rm f}\ge0$ with
$a_{\rm r}+a_{\rm f}\le q$, together with sup-norm majorants $e_j$ of
the regular atoms and $e'_a\ge0$ with $|k_a|\le e'_a$, such that
\[
 \sum_j\operatorname{mass}(X_j,e_j)\le a_{\rm r},\qquad
 \sum_a\operatorname{mass}(S_a,e'_a)\le\kappa\,a_{\rm f}.
\]
Here $\operatorname{mass}(X,e)=e^{\alpha_{i+1}\operatorname{coverNumber}(B_{i+1},X)}e$
is the fine mass. It is linear in $e$ for a fixed support. This is the
normalized remainder convention of Section~\ref{sec:local-splitting}.

A \emph{Coupon} across $D$ on a wing $W$ for input data
$(\alpha,\rho,\delta,f)$ is defined as follows. Here $\alpha$ is an
exponential entry, $\rho$ a forced profile, $\delta$ a unit value and
$f$ an upper density. The Coupon has lead $\ell_0\in F$ and prices
$(c,r,u)$. It is a lower coefficient field $P$ on $R^-$ with the
format at $R^-$, the exact push identity
$\operatorname{dens}(P)=L_i^d\mathsf P_if$, the lead bound
$C_{\rm orb}\|P.x-\ell_0\|_1\le c$, and a turn-weighted remainder
bound at most $r$. It also has a Tale witness on $W$ across $D$
relating $(\alpha,\rho,\delta)$ to the lower forced profile
$\mathcal P(P)$, with input price $u$ and output price $r$. This is the
notion of Section~\ref{sec:three-port-history}: there the three
elementary insertions are described by exactly these items, with
prices written as $C_{\rm orb}\|Q.x-\cdots\|_1\le\cdot$,
$r=\cdot$ and ``input allowance''. It is also the description before
Lemma~\ref{lem:actual-marked-antilles}, with $\ell_0=0$.

\begin{proposition}[Zero-coordinate Coupon on the stored doors]
\label{prop:selected-zero-coordinate-coupon}
Keep the selected history, its stored doors and the purchase of
Section~\ref{sec:common-observable-purchases}. Fix $d,N,p$ and
$t_d\le i<N$, with the notation above. Assume the following three
inputs on the stored door $D$.
\begin{enumerate}
\item[(K)] \emph{Insertion estimates.} The upper representation $R^+$
  has a Keeper, an injective assignment $(X,I)\mapsto\mathsf k(X,I)$
  of an address with tag $(X,I)$ to every pair
  $\varnothing\ne I\subseteq X$, eligible whenever
  $X\cap\operatorname{flags}(U)=I$. Its keys lie in $W_i$. On the wing
  $W_i$ the regular and forced estimates of
  Section~\ref{sec:three-port-history} hold on $D$ with $\eta=\varrho_i$,
  in the following form, with lead $0$.
  \begin{itemize}
  \item Let $(X,f)$ be a regular atom with $0\in X$ that passes the
    format tests, and let $e$ be a sup-norm majorant of $f$. Put
    $m=\operatorname{mass}(X,e)$. Then the input data
    $(\chi_Xf,\,0,\,0,\,R^+.\mathrm{eval}\,\chi_Xf)$ have a Coupon with
    prices $(\varrho_im,\varrho_im,2m)$.
  \item Let $a$ be an address whose key lies in $W_i$, and let
    $f\ne0$ be a forced insertion at $a$ that fits its slot, with
    $|f|\le e$. Put $m=\operatorname{mass}(S_a,e)$. Then the input data
    $(0,\,f\mathbf 1_{\operatorname{tag}(a)},\,0,\,
    e^{R^+.\mathrm{exponent}}f\operatorname{Dr}_a)$ have a Coupon with
    prices $(\varrho_im/\kappa,\varrho_im/\kappa,2m)$.
  \end{itemize}
\item[(S)] \emph{Soundness.} The door $D$ is sound for the transition
  operator of the step.
\item[(M)] \emph{Merger.} Coupons across $D$ on one wing are monotone
  in their three prices. A countable family of them merges into a
  Coupon whose input data, lead and prices are the sums, provided the
  following hold: the entries and the profiles converge pointwise; the
  unit values, the three prices and the norms of the leads are
  summable; and the densities are measurable and dominated. This is the
  merger left as an outstanding input in
  Section~\ref{ex-selected-linear-interpolation}, together with the
  summation of the Tale witnesses. Finally, zero input data have a
  Coupon with lead $0$ and any nonnegative prices.
\end{enumerate}
Then (R7) holds at $D$ with $\eta_i=\varrho_i$:
\[
 (1+2C_i)\varrho_i\le\operatorname{barunit}(L_i),\qquad
 D\text{ is sound},\qquad \psi_i\ge0,\qquad \varrho_i>0.
\]
Moreover, let $Q$ be any upper coefficient field on $R^+$ with
$Q.x=0$, the terminal format at the origin and remainder bound $q$ at
weight $\kappa$. Then its input data
$(\mathcal E(Q),\mathcal P(Q),Q.x_0,\operatorname{dens}(Q))$ have a
Coupon across $D$ on $W_i$ with lead $0$ and prices
$(\varrho_iq,\varrho_iq,2q)$. The wing $W_i$ is finite and does not
depend on $Q$.
\end{proposition}

\begin{proof}
\emph{Tuning.} Section~\ref{sec:ex-marked-portfolio-calibration}
chooses, for every stage $j$, a threshold $v_j$ with
$(1+2c_j)\operatorname{rate}(L)\le\operatorname{barunit}(L)$ for all
$L\ge v_j$; here $c_j$ is the column count of the profile at stage $j$.
Indeed the ratio of the two sides is
$\operatorname{rate}(L)/\operatorname{barunit}(L)=L^{-2/5+3/8}=L^{-1/40}$.
So the inequality is equivalent to $L^{1/40}\ge1+2c_j$, and
$v_j=\lceil(1+2c_j)^{40}\rceil$ suffices. The stage threshold $n_k$
there is the maximum of these $v_j$ over $j\le k$, together with the
credit thresholds. By Section~\ref{sec:common-observable-purchases},
the purchase includes $n_k$ among the block floors of stage $k$. Every
block length chosen at stage $k$ is at least that floor. Take
$k=J(i+1)$. The block side of the present step is chosen at that stage,
so $L_i\ge n_{J(i+1)}$. Since $J(i)\le J(i+1)$, the threshold
$v_{J(i)}$ is one of the terms of $n_{J(i+1)}$. Hence
$L_i\ge v_{J(i)}$. The column count at level $i$ is that of the profile
at stage $J(i)$, so $C_i=c_{J(i)}$. Therefore
$(1+2C_i)\varrho_i\le\operatorname{barunit}(L_i)$. This is the
grade-zero clause of the purchase: the inequality involves no section,
and the threshold $v_j$ is the same for every grade.

\emph{Signs and soundness.} We have $\varrho_i=L_i^{-2/5}>0$. Also
$\psi_i=2e^{\alpha_{i+1}}\tau_{i+1}\ge0$, since $\tau_{i+1}\ge0$.
Soundness is (S).

\emph{Repacking the forced array.} Let $Q$ be as in the statement. Define
$Q'$ to have the same coordinate vector and regular atoms as $Q$. Its
forced array is
\[
 k'_{\mathsf k(X,I)}=\mathcal P(Q)(X,I),\qquad
 k'_a=0\ \text{ if $a$ is not a Keeper address}.
\]
The tag of every address has a nonempty inventory contained in its
support, so it lies in the domain of the Keeper. We check five
properties of $Q'$.

(i) $\mathcal E(Q')=\mathcal E(Q)$, because the exponential entry does
not involve the forced array.

(ii) $\mathcal P(Q')=\mathcal P(Q)$. Fix a tag $(X,I)$. The only address
with this tag at which $k'$ can be nonzero is $\mathsf k(X,I)$, since
$\mathsf k$ is injective and $\operatorname{tag}(\mathsf k(X,I))=(X,I)$.
Hence $\mathcal P(Q')(X,I)=k'_{\mathsf k(X,I)}=\mathcal P(Q)(X,I)$.

(iii) $\operatorname{dens}(Q')=\operatorname{dens}(Q)$. The dressed
field depends on the address only through its tag. To see this, let
$\operatorname{tag}(a)=\operatorname{tag}(a')$ with $a\ne a'$, and let
$\sigma$ be the transposition of $a$ and $a'$. Then $\sigma$ preserves
tags, so it maps admissible sets to admissible sets. No admissible set
contains both $a$ and $a'$, because their supports coincide and are
nonempty, while distinct members of an admissible set have disjoint
supports. Thus $b\mapsto\sigma(b)$ is a bijection from the admissible
sets containing $a$ to those containing $a'$, and
$\sigma(b)\setminus\{a'\}=b\setminus\{a\}$. Comparing the two series
term by term gives $\operatorname{Dr}_a=\operatorname{Dr}_{a'}$. Write
$\operatorname{Dr}_{(X,I)}$ for the common value. The forced series
converges absolutely. Grouping it by tags therefore gives
\[
 \sum_ak_a(U)\operatorname{Dr}_a(U)
 =\sum_{(X,I)}\mathcal P(Q)(X,I)(U)\operatorname{Dr}_{(X,I)}(U),
\]
and the same identity holds for $Q'$. By (i) and (ii) the two
densities coincide.

(iv) $Q'$ has the terminal format. The regular clauses are unchanged.
Fix a tag $(X,I)$ and put $g=k'_{\mathsf k(X,I)}=\sum_{\operatorname{tag}(a)=(X,I)}k_a$.
This sum converges absolutely, by the majorants $e'_a$ and the
summability below. A countable absolutely convergent sum of
gauge-invariant, measurable functions local to $X$ has the same three
properties. If $g(U)\ne0$, then some $k_a(U)\ne0$ with
$\operatorname{tag}(a)=(X,I)$. Eligibility of $a$ at $U$ gives
$X\cap\operatorname{flags}(U)=I$, so the Keeper address $\mathsf k(X,I)$
is eligible at $U$. Similarly $g\ne0$ forces some $k_a\ne0$, hence
$0\in S_a=X$.

(v) $Q'$ has remainder bound $q$ at weight $\kappa$, with the same
$a_{\rm r},a_{\rm f}$. The regular part is unchanged. For the forced
part, put $e''_{\mathsf k(X,I)}=\sum_{\operatorname{tag}(a)=(X,I)}e'_a$
and $e''=0$ elsewhere. Then $|k'|\le e''$ by the triangle inequality.
Masses on the common support $X$ add, so
$\sum_a\operatorname{mass}(S_a,e''_a)=\sum_a\operatorname{mass}(S_a,e'_a)\le\kappa a_{\rm f}$.
The mass bound gives $\operatorname{mass}(S_a,e'_a)\ge e'_a\ge0$
(the cover number is nonnegative), so $\sum_ae'_a<\infty$. This is the
summability used in (iv).

Consequently $Q$ and $Q'$ have the same input data
$(\mathcal E,\mathcal P,Q.x_0,\operatorname{dens})$, and it suffices to
produce the Coupon for $Q'$.

\emph{Decomposition.} Since $Q'.x=0$, also $Q'.x_0=0$, and the placement
term $\ell(U,0)$ vanishes. Thus
\[
 \mathcal E(Q')=\sum_j\chi_{X_j}f_j,\qquad
 \mathcal P(Q')=\sum_a k'_a\mathbf 1_{\operatorname{tag}(a)},\qquad
 \operatorname{dens}(Q')=\sum_jR^+.\mathrm{eval}\,\chi_{X_j}f_j
   +\sum_ae^{R^+.\mathrm{exponent}}k'_a\operatorname{Dr}_a .
\]
All three series converge absolutely and pointwise. For the density the
regular and forced majorants are dominated by the coefficient envelope
of the lawful representation and by the remainder bound. Tonelli's
theorem applied to these nonnegative majorants gives a dominated,
measurable expansion, as in
Section~\ref{ex-selected-linear-interpolation}. Put
$m_j=\operatorname{mass}(X_j,e_j)$ and
$m'_a=\operatorname{mass}(S_a,e''_a)$. We attach a Coupon to each term.
\begin{itemize}
\item Regular atom $j$. Its input data are
  $(\chi_{X_j}f_j,0,0,R^+.\mathrm{eval}\,\chi_{X_j}f_j)$. The atom has
  $0\in X_j$, passes the format tests and has majorant $e_j$. By (K) it
  has a Coupon with prices $(\varrho_im_j,\varrho_im_j,2m_j)$.
\item Forced address $a$ with $k'_a\ne0$. Then $a=\mathsf k(X,I)$ is a
  Keeper address, so its key lies in $W_i$. Its coefficient fits the
  slot by (iv), and $|k'_a|\le e''_a$. By (K) it has a Coupon with
  prices $(\varrho_im'_a/\kappa,\varrho_im'_a/\kappa,2m'_a)$. Since
  $\kappa\le1$, monotonicity (M) gives one with prices
  $(\varrho_im'_a/\kappa,\varrho_im'_a/\kappa,2m'_a/\kappa)$.
\item Forced address $a$ with $k'_a=0$. Its input data are zero, and
  by (M) they have a Coupon with prices
  $(\varrho_im'_a/\kappa,\varrho_im'_a/\kappa,2m'_a/\kappa)$.
\end{itemize}

\emph{Summation.} All leads are $0$, and so are all unit values. The
entries, profiles and densities of the terms sum to those of $Q'$, by
the decomposition. The price series are summable:
\[
 \sum_j\varrho_im_j+\sum_a\frac{\varrho_im'_a}{\kappa}
 \le\varrho_i\bigl(a_{\rm r}+a_{\rm f}\bigr)\le\varrho_iq,
 \qquad
 \sum_j2m_j+\sum_a\frac{2m'_a}{\kappa}\le2q .
\]
Here we used $\sum_jm_j\le a_{\rm r}$ and
$\sum_am'_a\le\kappa a_{\rm f}$ from (v). By (M) the terms merge into
one Coupon on $W_i$ for the input data of $Q'$, with lead $0$. Its
prices are the three sums, and by monotonicity they may be replaced by
$(\varrho_iq,\varrho_iq,2q)$. By the repacking step these are also the
input data of $Q$.

The wing $W_i$ was chosen before $Q$, from the finite fine torus and
its Keeper. It is finite because the torus has finitely many supports
and tags.
\end{proof}

\paragraph{What the proposition supplies.}
Of (R7), the proposition proves the tuning inequality, the signs, and
the reduction of the Coupon for an arbitrary zero-coordinate
coefficient to single insertions. The tuning inequality uses only the
grade-zero threshold of the purchase. That threshold pays tuning; the
Coupon itself is paid at the grade $d$ of the frame, through (K). The
three inputs are not proved here.
\begin{itemize}
\item The regular and forced estimates (K) are those of
  Section~\ref{sec:three-port-history}. They hold relative to the
  retained funded interpolation and its slot realization, with
  $\eta=L^{-2/5}$ there equal to $\varrho_i$. They are attached to the
  stored door through the same history identifications. Their
  derivation from the selected interpolation is the open path
  construction recorded at the end of that section and in
  Section~\ref{ex-centered-freight-producer}.
\item Soundness (S) of the stored door is asserted in
  Section~\ref{sec:three-port-history} for the retained interpolating
  point.
\item The merger (M) is the outstanding input of
  Section~\ref{ex-selected-linear-interpolation}, extended to the Tale
  witnesses.
\end{itemize}
The Keeper of $R^+$ in (K) is part of the retained slot realization.
At the terminal level it is the Keeper of
Proposition~\ref{prop:terminal-marked-voice}.
